\documentclass[11pt,a4paper]{article}
\usepackage[utf8]{inputenc}
\usepackage[T1]{fontenc}
\usepackage[brazil]{babel}
\usepackage{lmodern}
\usepackage{microtype}
\usepackage{geometry}
\geometry{margin=2.25cm}
\usepackage{amsmath,amssymb,booktabs,longtable,array,xcolor,graphicx,hyperref}
\hypersetup{colorlinks=true,linkcolor=blue!50!black,citecolor=blue!50!black,urlcolor=blue!60!black}
\setlength{\parindent}{0pt}
\setlength{\parskip}{0.55em}
\newcommand{\code}[1]{\texttt{#1}}
\newcommand{\rhoNine}{\rho_9}

\title{Validação de PYTHIA 8 e GEANT4 em matéria densa no HADROS3\\
\large Revisão bibliográfica, limites físicos e campanha quantitativa}
\author{Projeto HADROS3}
\date{14 de julho de 2026}

\begin{document}
\maketitle

\begin{abstract}
Este relatório confronta a cadeia HADROS3 de interação neutrino--núcleon,
geração do estado final e transporte com resultados modernos de QCD,
medições de neutrinos no LHC, testes de calorímetros e teoria de matéria
densa. A conclusão central é que há três problemas fisicamente distintos:
(i) a probabilidade e a cinemática da interação DIS; (ii) o shower e a
hadronização de uma colisão isolada; e (iii) a propagação coletiva dos
secundários no plasma denso. PYTHIA 8 é apropriado para (ii), e GEANT4 possui
forte validação laboratorial para partes de (iii), mas nenhum experimento
terrestre reproduz o plasma H/He degenerado de
$\rho\sim10^9$--$10^{10}\,\mathrm{g\,cm^{-3}}$. Portanto, a execução atual é
uma extrapolação controlável, não uma validação completa do meio
astrofísico. São definidos testes numéricos, observáveis e critérios de
aceitação para tornar essa afirmação quantitativa.
\end{abstract}

\tableofcontents

\section{Pergunta física e escopo}

O objetivo é saber quanto se pode confiar na produção e no transporte de
partículas nos pontos do toro amostrados pelo HADROS3. É essencial não
misturar densidade macroscópica com física do vértice elementar:
{\small\begin{equation}
\underbrace{\nu N\rightarrow \ell/X}_{\text{DIS e PDFs}}
\quad\longrightarrow\quad
\underbrace{\text{shower, hadronização, decaimentos}}_{\text{PYTHIA}}
\quad\longrightarrow\quad
\underbrace{\text{propagação e reinterações no meio}}_{\text{GEANT4 + correções densas}}.
\end{equation}}
Em primeira aproximação, a densidade local multiplica a taxa de alvos e
fixa profundidades ópticas; ela não deve alterar diretamente uma colisão
isolada no PYTHIA. Efeitos do meio entram nas PDFs nucleares, no estado
termodinâmico dos alvos e, sobretudo, na propagação dos secundários.

\section{Estado concreto do HADROS3 analisado}

Os produtos locais de H3-W10 registram PYTHIA 8.312 e HepMC3 3.03.01, dez
eventos, 298 partículas finais e 145 hádrons finais. A maior violação
relativa de conservação de quatro-momento é $3.71\times10^{-10}$, bem abaixo
do critério $5\times10^{-8}$. A taxa medida é $4.77$ eventos/s.

H3-W11 transportou os dez eventos em dez processos isolados por sítio,
usando GEANT4 11.4.2, a lista FTFP\_BERT e um material H/He homogêneo com
frações de massa $X=0.75$ e $Y=0.25$. Cada caixa tem semi-lado de 10 mm. As
densidades vêm dos pontos DIS e alcançam aproximadamente
$1.36\times10^9$ a $9.81\times10^9\,\mathrm{g\,cm^{-3}}$ na amostra exibida.
A execução importou 298 primárias, alcançou energia máxima de
$9.45\times10^3$ GeV, efetuou $2.734\times10^7$ passos em 43.1 s e truncou o
registro em $5\times10^5$ passos. Esses números demonstram execução e
consistência interna; não demonstram por si sós validade do modelo de meio.

\section{O que os papers permitem validar}

\subsection{PYTHIA e o evento elementar}

O manual moderno do PYTHIA 8.3 descreve shower, hadronização por strings,
decaimentos, tunes e interfaces LHA/HepMC \cite{Bierlich2022Pythia}. Para o
HADROS3, a comparação apropriada é no mesmo evento POWHEG/LHE: cardinalidade
1:1, conservação de quatro-momento, conteúdo de espécies, multiplicidade,
escala de matching e estabilidade sob seeds. A densidade do toro deve ficar
como metadado do sítio, nunca como um ``tune de densidade'' do PYTHIA.

O resultado atual de conservação é excelente. Ainda faltam testes de forma:
distribuições de multiplicidade, fração de energia por espécie, $p_T$,
rapidez e dependência com a inelasticidade $y$. Esses testes devem comparar
PYTHIA e, quando possível, uma segunda cadeia geradora ou dados de neutrinos.

\subsection{DIS em TeV e UHE}

Xie et al. calculam seções de choque de DIS de alta energia com PDFs modernas
e incertezas \cite{Xie2024}. Weigel et al. fornecem seções totais e
distribuições de inelasticidade de 50 GeV a $5\times10^{12}$ GeV, incluindo
massas de quarks pesados, small-$x$, correções nucleares e radiação de estado
final \cite{Weigel2024}. Le e Mäntysaari quantificam supressão nuclear e a
componente difrativa em neutrino--núcleo \cite{Le2024}.

Esses trabalhos são o benchmark direto para a camada DIS do HADROS3. Em
cada bin de $E_\nu$, devem ser comparados $\sigma_{CC}$, $\sigma_{NC}$,
$d\sigma/dy$, razão $\bar\nu/\nu$ e incertezas. Como o toro é H/He, a
correção nuclear deve ser aplicada por espécie e fração, não pela densidade
de massa. A densidade determina $n_N\sigma$, enquanto estrutura nuclear e
small-$x$ determinam $\sigma$.

\subsection{O elo experimental do CERN}

O FASER mediu diretamente $\nu_e$ e $\nu_\mu$ em alvo de tungstênio, com
energias reconstruídas aproximadamente entre 0.5 e 1.8 TeV
\cite{FASER2024}. A análise posterior mediu a seção de choque de
$\nu_\mu+\bar\nu_\mu$ diferencial em energia, em seis bins, usando 338
eventos CC \cite{FASER2025}. É o teste experimental mais próximo da escala
TeV do HADROS3 entre os papers selecionados.

Esse dado valida o setor eletrofraco/DIS e partes da modelagem de estados
finais, mas não a densidade astrofísica: tungstênio de detector é matéria
fria ordinária. Além disso, a energia atual máxima do HADROS3, 9.45 TeV,
fica acima do alcance direto publicado pelo FASER, exigindo extrapolação com
PDFs e sua incerteza.

\subsection{GEANT4, FTFP\_BERT e test beams}

GEANT4 é uma biblioteca geral de transporte cujo desenvolvimento e domínio
de aplicações são descritos por Allison et al. \cite{Allison2016Geant4}.
CALICE confrontou modelos hadrônicos GEANT4 com chuveiros em calorímetros de
aço altamente granulares \cite{CALICE2013,CALICE2016,CALICE2019}. Esses
papers permitem testar perfis longitudinais e radiais, energia visível,
tempo, multiplicidade de hits e decomposição em componentes eletromagnética,
hadronicamente interagente, nêutrons e fragmentos.

Eles sustentam o uso de FTFP\_BERT como referência para chuveiros em materiais
laboratoriais e nas faixas de energia testadas. Não sustentam extrapolar um
\code{G4Material} atômico para $10^9$--$10^{10}\,\mathrm{g\,cm^{-3}}$ sem
correções. O teste CALICE deve ser reproduzido separadamente no HADROS3 como
benchmark do executável; não se deve ajustar o toro para ``parecer'' um
calorímetro.

\section{Por que a densidade extrema muda o problema}

Para $X=0.75$, $Y=0.25$, a fração eletrônica é
$Y_e=X+Y/2=0.875$. Definindo $\rhoNine=\rho/(10^9\,\mathrm{g\,cm^{-3}})$,
\begin{align}
n_e &\simeq 5.27\times10^{32}\rhoNine\ \mathrm{cm^{-3}},\\
p_Fc &= \hbar c(3\pi^2n_e)^{1/3}\simeq4.93\rhoNine^{1/3}\ \mathrm{MeV},\\
E_{F,\mathrm{kin}}&=\sqrt{(p_Fc)^2+(m_ec^2)^2}-m_ec^2.
\end{align}
Assim, $E_{F,\mathrm{kin}}\simeq4.45$ MeV em $10^9$ e cerca de 10.1 MeV em
$10^{10}\,\mathrm{g\,cm^{-3}}$. Elétrons são relativisticamente degenerados
se $kT\ll E_F$. A densidade de íons é aproximadamente
$4.89\times10^{32}\rhoNine\,\mathrm{cm^{-3}}$ e o raio de Wigner--Seitz é
aproximadamente $79\rhoNine^{-1/3}$ fm: muito menor que uma escala atômica,
mas ainda maior que raios nucleares.

Esse regime é conhecido na física de crostas de estrelas compactas e
transporte degenerado \cite{Chamel2008,Potekhin2015}. Ele implica ionização
completa, screening, correlações do plasma, ocupação de estados eletrônicos e
possível bloqueio de Pauli. Para processos eletromagnéticos de alta energia,
o comprimento de formação pode introduzir supressões LPM e dielétrica
\cite{Klein1999}. Estudos modernos de stopping em matéria densa mostram
explicitamente que campos coletivos podem tornar inadequada a soma de perdas
por colisões independentes \cite{Fu2023}. O paper de Fu et al. trata alfas em
DT de fusão, não hádrons TeV no toro; sua importância é demonstrar a classe
de física ausente, não fornecer uma correção pronta.

Há uma nuance importante: os núcleos ainda estão separados por dezenas de
fm, portanto uma colisão forte dura em um núcleo pode continuar sendo tratada
localmente. O maior risco não é que o PYTHIA ``veja todos os núcleos ao mesmo
tempo'', mas que o transporte entre colisões, a resposta eletromagnética e o
resfriamento/absorção de secundários sejam modelados como matéria atômica
fria independente.

\section{Fontes ocultas e jatos dentro de matéria}

Murase e Ioka mostram que jatos dentro de estrelas podem produzir neutrinos
TeV--PeV, mas que a radiação mediada e o resfriamento dos píons/múons impõem
condições fortes \cite{MuraseIoka2013}. Senno, Murase e Mészáros estendem a
análise para jatos sufocados e GRBs de baixa luminosidade, enfatizando fontes
ocultas e competição entre produção e perdas \cite{Senno2016}.

Para o HADROS3, a lição não é copiar a geometria estelar. É calcular, em cada
sítio, a hierarquia
\begin{equation}
t_{\rm dec},\quad t_{\rm had},\quad t_{\rm em},\quad t_{\rm syn},\quad
t_{\rm adv},
\end{equation}
para píons, káons, múons, nucleons e fótons. Se $t_{\rm had}<t_{\rm dec}$, a
composição final prevista apenas por decaimentos PYTHIA não representa o que
escapa do toro. O PYTHIA continua correto como condição inicial; a evolução
posterior é que muda.

\section{Escalas úteis e interpretação do volume atual}

Uma caixa de lado 20 mm atravessada integralmente tem coluna
\begin{equation}
X_{\rm box}=\rho L=2\times10^9\rhoNine\ \mathrm{g\,cm^{-2}}.
\end{equation}
Logo, os sítios atuais representam colunas de ordem $10^9$--$2\times10^{10}$
g cm$^{-2}$. Comparações de densidade devem manter a coluna $X=\rho L$, não
o comprimento geométrico, quando a aproximação de centros independentes é
válida. Para uma seção neutrino--núcleon $\sigma_{\nu N}$,
\begin{equation}
\lambda_X=(N_A\sigma_{\nu N})^{-1},\qquad
P_{\rm int}=1-\exp(-X/\lambda_X).
\end{equation}
Para hádrons, seções fortes são muitas ordens maiores que as fracas; por isso
os comprimentos físicos entre interações tornam-se microscópicos e o número
de passos cresce rapidamente. Os $2.7\times10^7$ passos atuais são coerentes
com um problema computacionalmente rígido, mas a rigidez não prova que as
seções e perdas usadas sejam as do plasma real.

\section{Matriz de validade}

{\footnotesize
\begin{longtable}{p{0.20\textwidth}p{0.24\textwidth}p{0.24\textwidth}p{0.17\textwidth}}
\toprule
Componente & Evidência disponível & Limitação & Status recomendado \\
\midrule\endhead
Conservação H3-W10 & testes numéricos do próprio run & não testa distribuições & validado internamente \\
Shower/hadronização & manual/tunes PYTHIA e dados HEP & canal DIS e energia específicos precisam comparação & parcialmente validado \\
Seção DIS & cálculos NNLO/NLO e FASER TeV & extrapolação até 10 TeV e small-$x$ & validável com incerteza \\
Chuveiro GEANT4 & CALICE aço, FTFP\_BERT & material/energia diferentes & validado apenas no benchmark \\
Escala $1/\rho$ & teoria de centros independentes & falha com efeitos coletivos & teste de consistência \\
Plasma H/He denso & teoria de matéria degenerada & sem dados de hádrons TeV nessas densidades & não validado \\
Toro completo & modelo analítico local & sem GRMHD, campos ou gradientes no patch & aproximação exploratória \\
\bottomrule
\end{longtable}
}

\section{Campanha numérica proposta}

\subsection*{V1 --- DIS contra referências modernas}
Gerar uma grade logarítmica de 50 GeV a $10^5$ GeV, com pelo menos $10^5$
amostras por década. Comparar $\sigma_{CC/NC}$ e quantis de $y$ com Xie e
Weigel. Aceitar se a curva central estiver dentro da banda teórica publicada;
onde não houver tabela digital, exigir acordo de 5\% após digitalização e
propagar a incerteza de PDF separadamente.

\subsection*{V2 --- fechamento do PYTHIA}
Em cada bin de energia e $y$, gerar pelo menos $10^4$ eventos. Exigir
cardinalidade 1:1, carga exata, resíduo relativo de quatro-momento abaixo de
$5\times10^{-8}$ em 100\% dos eventos, ausência de pártons finais e nenhuma
emissão acima de SCALUP. Publicar médias, covariâncias e intervalos de
multiplicidade/fração de energia por espécie. A densidade não pode alterar o
hash do evento PYTHIA para a mesma seed e LHE.

\subsection*{V3 --- benchmark GEANT4/CALICE}
Construir uma geometria de validação separada em aço com feixes de $\pi^\pm$
nas energias dos papers. Comparar energia média visível, resolução, centro de
gravidade longitudinal e raio do shower. Aceitar somente quando cada
observável estiver dentro da incerteza experimental combinada ou, na ausência
de covariância, dentro de 10\%. Esse teste valida executável/lista física, não
o toro.

\subsection*{V4 --- invariância em profundidade de massa}
Usar o mesmo material ideal diluído em uma região onde efeitos coletivos são
desprezíveis. Executar pares $(\rho,L)$ com $\rho L$ fixo e a mesma seed.
Comparar probabilidade da primeira interação, energia depositada por coluna e
espectro de escape. Exigir desvio menor que 2\% com pelo menos $10^4$ eventos.
Falha indica dependência geométrica/numérica indevida; sucesso não valida o
regime degenerado.

\subsection*{V5 --- convergência de produção cuts e passos}
Variar production cut e step limiter por fatores 2 e 4. Aceitar quando energia
depositada, energia escapada e multiplicidades relevantes variarem menos que
1\%, enquanto o balanço e os hashes de entrada permanecem válidos. O atual
\code{steps\_truncated=true} deve ser reportado como limitação de saída, nunca
como perda silenciosa de física.

\subsection*{V6 --- envelope de matéria densa}
Calcular $E_F$, parâmetro de acoplamento iônico, comprimentos de screening,
formação e interação em cada sítio usando também a temperatura local. Rodar
três modelos: GEANT4 atômico bruto; stopping com correções de plasma; e
limites superior/inferior para LPM, screening e bloqueio. Não há critério de
``acordo com dado''; o produto aceito é uma banda sistemática. Se a banda
exceder 20\% num observável científico, ele deve ser marcado como dependente
do modelo de meio.

\subsection*{V7 --- competição interação/decaimento}
Para $\pi^\pm$, $K^\pm$, $K_L$, $\mu^\pm$, nêutrons e fótons, comparar
$t_{\rm dec}$, $t_{\rm int}$, $t_{\rm loss}$ e $t_{\rm escape}$ por energia e
sítio. Exigir que toda partícula declarada ``final'' tenha a decisão
transportar/decaír/absorver documentada. Este teste liga diretamente PYTHIA
ao observável astrofísico.

\subsection*{V8 --- desempenho}
Medir tempo por evento, passos por GeV e memória em função de densidade,
coluna e multiplicidade. Com o run atual, a referência é 10 eventos em 43.1 s
com quatro workers. Definir um orçamento de erro, não apenas de tempo:
parallelismo por sítio é seguro; juntar densidades em um mesmo material não é.
Aplicar thinning/biasing somente com pesos auditados e comparação contra uma
amostra sem aceleração, exigindo acordo estatístico de 2\%.

\section{Mudanças recomendadas no HADROS3}

\begin{enumerate}
\item Separar no relatório web três selos independentes: validação do gerador
de eventos, validação do benchmark GEANT4 e validação da microfísica do meio
denso. O terceiro deve começar falso; os nomes de máquina propostos são
\code{event\_generator\_validated}, \code{geant4\_benchmark\_validated} e
\code{dense\_medium\_microphysics\_validated}.
\item Registrar temperatura, ionização, $Y_e$, degenerescência, magnetização e
gradientes junto à densidade. Densidade sozinha não define um material.
\item Reportar comprimentos e histogramas também em g cm$^{-2}$.
\item Preservar o processamento por sítio já implementado: um material/volume
local para cada densidade. Isso corrige a mistura de sítios, mas não corrige a
microfísica de plasma.
\item Criar um backend de correções densas versionado, inicialmente como
envelope sistemático, sem modificar PYTHIA.
\item Adicionar os testes V1--V8 à campanha de física e gravar DOI/arXiv,
versões, seeds, cards e hashes em cada produto.
\end{enumerate}

\section{Conclusão}

É possível validar grande parte da cadeia, mas não com uma única comparação.
FASER e os cálculos modernos testam a interação de neutrinos em TeV; PYTHIA
testa o estado final de cada colisão; CALICE testa FTFP\_BERT em calorímetros.
Os papers de matéria densa mostram por que a etapa restante exige nova
microfísica. O resultado atual do HADROS3 deve ser descrito como
\emph{transporte GEANT4 de uma mistura H/He homogênea escalada à densidade
local, com consistência interna verificada e efeitos coletivos ainda não
validados}. Essa formulação é rigorosa, útil e aponta diretamente para os
testes que podem convertê-la, progressivamente, em uma previsão astrofísica.

\bibliographystyle{unsrt}
\bibliography{references}

\end{document}
